\documentclass{article}

\usepackage{fancyhdr}
\usepackage{extramarks}
\usepackage{amsmath}
\usepackage{amsthm}
\usepackage{amsfonts}
\usepackage{tikz}
\usepackage[plain]{algorithm}
\usepackage{algpseudocode}
\usepackage{fullpage, url, hyperref}


\theoremstyle{plain}
\newtheorem{theorem}{Theorem}
\newtheorem{lemma}{Lemma}
\newtheorem{definition}{Definition}

\newcommand{\I}{\mathcal{I}}
\newcommand{\F}{\mathbb{F}}
\newcommand{\E}{\mathbb{E}}
\newcommand{\R}{\mathbb{R}}
\newcommand{\tr}{\mathrm{tr}}
\renewcommand{\P}{\mathbb{P}}
\usetikzlibrary{automata,positioning}

%
% Basic Document Settings
%

\topmargin=-0.45in
\evensidemargin=-.50in
\oddsidemargin=-.5in
\textwidth=7in
\textheight=9.5in
\headsep=0.25in

\linespread{1.1}

\pagestyle{fancy}
\lhead{\hmwkAuthorName}
\chead{\hmwkClass\ (\hmwkClassInstructor\ \hmwkClassTime): \hmwkTitle}
\rhead{\firstxmark}
\lfoot{\lastxmark}
\cfoot{\thepage}

\newcommand{\w}[2]{w_{#1}^{(#2)}}
\newcommand{\loss}[2]{\vec{\ell}_{#2}^{(#1)}}
\newcommand{\losst}[1]{\vec{\ell}^{(#1)}}
  \newcommand{\pt}[1]{\vec{p}^{(#1)}}
\renewcommand\headrulewidth{0.4pt}
\renewcommand\footrulewidth{0.4pt}

\setlength\parindent{0pt}


% 
% Create Problem Sections
%


\newcommand{\enterProblemHeader}[1]{
    \nobreak\extramarks{}{}\nobreak{}
    \nobreak\extramarks{}\nobreak{}
}

\newcommand{\exitProblemHeader}[1]{
    \nobreak\extramarks{}\nobreak{}
    \stepcounter{#1}
    \nobreak\extramarks{}\nobreak{}
  }

  
\newcommand{\enterExerciseHeader}[1]{
    \nobreak\extramarks{}{}\nobreak{}
    \nobreak\extramarks{}\nobreak{}
}

\newcommand{\exitExerciseHeader}[1]{
    \nobreak\extramarks{}\nobreak{}
    \stepcounter{#1}
    \nobreak\extramarks{}\nobreak{}
}

\setcounter{secnumdepth}{0}
\newcounter{partCounter}
\newcounter{homeworkProblemCounter}
\setcounter{homeworkProblemCounter}{1}

\newcounter{homeworkExerciseCounter}
\setcounter{homeworkExerciseCounter}{1}
\nobreak\extramarks{}\nobreak{}
\nobreak\extramarks{}\nobreak{}

%
% Homework Problem Environment
%
% This environment takes an optional argument. When given, it will adjust the
% problem counter. This is useful for when the problems given for your
% assignment aren't sequential. See the last 3 problems of this template for an
% example.
%
\newenvironment{homeworkProblem}[1][-1]{
%    \ifnum#1>0
   %     \setcounter{homeworkProblemCounter}{#1}
   % \fi
    \section{Problem \arabic{homeworkProblemCounter}}
    \setcounter{partCounter}{1}
    \enterProblemHeader{homeworkProblemCounter}
}{
    \exitProblemHeader{homeworkProblemCounter}
}

  \newenvironment{homeworkExercise}[1][-1]{
    \ifnum#1>0
        \setcounter{homeworkExerciseCounter}{#1}
    \fi
    \section{Exercise \arabic{homeworkExerciseCounter}}
    \setcounter{partCounter}{1}
    \enterExerciseHeader{homeworkExerciseCounter}
}{
    \exitExerciseHeader{homeworkExerciseCounter}
}

%
% Homework Details
%   - Title
%   - Due date
%   - Class
%   - Section/Time
%   - Instructor
%   - Author
%

\newcommand{\hmwkTitle}{Homework\ \#4}
\newcommand{\hmwkAssignedDate}{March 14}
\newcommand{\hmwkDueDate}{March 30, midnight}
\newcommand{\hmwkClass}{Design and Analysis of Algorithms, CS 6550}
\newcommand{\hmwkClassTime}{}%Section A}
\newcommand{\hmwkClassInstructor}{Professor Jamie Morgenstern}
\newcommand{\hmwkAuthorName}{}%\textbf{Josh Davis} \and \textbf{Davis Josh}}

%
% Title Page
%

\title{    \vspace{2in}
    \textmd{\textbf{\hmwkClass:\ \hmwkTitle}}\\
    \normalsize\vspace{0.1in}\small{Due\ on\ \hmwkDueDate\ at 3:10pm}\\
    \vspace{0.1in}\large{\textit{\hmwkClassInstructor\ \hmwkClassTime}}
    \vspace{3in}
}

\author{\hmwkAuthorName}
\date{}

\renewcommand{\part}[1]{\textbf{\large Part \Alph{partCounter}}\stepcounter{partCounter}\\}

%
% Various Helper Commands
%

% Useful for algorithms
\newcommand{\alg}[1]{\textsc{\bfseries \footnotesize #1}}

% For derivatives
\newcommand{\deriv}[1]{\frac{\mathrm{d}}{\mathrm{d}x} (#1)}

% For partial derivatives
\newcommand{\pderiv}[2]{\frac{\partial}{\partial #1} (#2)}

% Integral dx
\newcommand{\dx}{\mathrm{d}x}

% Alias for the Solution section header
\newcommand{\solution}{\textbf{\large Solution}}

% Probability commands: Expectation, Variance, Covariance, Bias
\newcommand{\Var}{\mathrm{Var}}
\newcommand{\Cov}{\mathrm{Cov}}
\newcommand{\Bias}{\mathrm{Bias}}
\usepackage{enumitem}
\newcommand\litem[1]{\item{\bfseries#1.\space}}

\newcommand{\Z}{\mathbb{Z}}

\begin{document}

% \maketitle

\pagebreak

{\bf Homework Out: \hmwkAssignedDate}\\
{\bf Due Date:  \hmwkDueDate}\\

The HW contains some exercises (fairly simple problems to check you
are on board with the concepts; don’t submit your solutions), and
problems (for which you should submit your solutions, and which will
be graded). Some problems have sub-parts that are exercises.  For this
problem set, it’s OK to work with others. (Groups of 2, maybe 3 max.)
That being said, please think about the problems yourself before
talking to others. Please cite all sources you use, and people you
work with. The expectation is that you try and solve these problems
yourself, rather than looking online explicitly for
answers. Submissions due at beginning of class on the due date.
Please check the Piazza for details on submitting your \emph{LaTeXed}
solutions.

\newcommand{\poly}{\textrm{poly}}
\paragraph{Problems}


\begin{enumerate}
  \litem{(A Counter, and the Median-of-Means Estimator.)} Here is a
  way of maintaining an approximate counter. (Call this the basic
  counter.)  Start with $X \gets 0$. When an element arrives, increment $X$ by
  $1$ with probability $2^{-X}$. When queried, return $N := 2^X − 1$.
  \begin{enumerate}
  \item Suppose the actual count is $n$, show that $\E[N] = n$, and
    $\textrm{Var}(N) = \frac{n(n−1)}{ 2}$ .

    Since its variance is large, average $k$ independent basic
    counters $N_1, N_2, \ldots, N_k$, and output the sample average
    $\hat{N} := \frac{1}{k}\sum_i N_i$ . Call this the $k$-mean counter.
  \item (Do not submit) Show that
   \[\P[\hat{N} \notin  (1 \pm \epsilon)n] \leq \frac{1}{2\epsilon^2 k}.\]
   Hence using $k = \frac{1}{P 2\epsilon^2 \delta}$ counters can make
     the failure probability at most $\delta$. (Said in other words,
     your error is less than $\epsilon n$ with “confidence”
     $1 -\delta$.) Here’s a way to use only
     $K = O( \frac{1}{ \epsilon^2}\log\frac{ 1}{\delta})$ counters to
     get the same answer (and the approach is useful in many different
     contexts beyond this one):

     Take a collection of $\ell = 10 log\frac{ 1 }{\delta}$
     independent $k_0$-mean counters, where $k_0 =\frac{4}{ \epsilon^ 2}$ .  Output the
     median $M$ of these $\ell$ counters.

   \item Prove that $\P[M \notin (1 \pm \epsilon ε)n] \leq
     \delta$. {\it Hint: what must happen for the the median to be too
       high?  What is the chance of that?}
   \end{enumerate}

   \litem{(I Stream, You Stream.)} In data streaming model, suppose we
   denote frequency vector by
   $x = (x_1, x_2, . . . , x_D) \in \Z^D_{\geq 0}$ where $x_i$ counts
   the number of occurences of element $i \in [D]$ seen so far. We
   want a streaming algorithm that stores information about the stream
   so that when it is eventually queried with some index $q \in [D]$,
   returns a value $\hat{x}_q$ that is $\approx x_q$ with probability
   at least $1 -\delta$. We don’t want to store $x$ explicitly, we want
   to use less space.

   Consider the following algorithm:

   Keep a global hash function $g : [D] \to [d]$, and also $d$ counters $C_1, C_2,\ldots, C_d$
(initially zero), each with its own hash function $h_i: [D] \to \{−1, +1\}$. If you see
element $e \in [D]$, first hash it using the global hash function $g$ to get the bucket
number $g(e)$, and then update

\[C_{g(e)}\gets   C_{g(e)} + h_{g(e)}(e).\]
When faced with the query $q$, output $A(q) := h_{g(q)}(q) · C_{g(q)}$.

Assume that $g$ and the $h_i$s are all independently picked, and each hash function is itself
$2$-universal.
\begin{enumerate}
\item  Show that $E[A(q)] = x_q$.
\item  Show that the variance of the estimate $A(q) = \frac{1}{d}(F_2 − x^2_q) \leq F^2/d.$
\item  Show that if we set $d =\frac{1}{\epsilon^2 \delta}$, our estimate $A(q)$ satisfies
\[\P[A(q) \in x_q \pm \epsilon \sqrt{F_2}] \geq 1 -\delta.\]
Recall that $F_2 := \sum_i x^2_i = || x||^2_2$  and hence the error term is $\epsilon||x||_2$.

\item Finally, consider an extension of this idea: maintain $t$
  independent copies of the above data structure. On a query for $q$,
  if the answers from the individual copies of the data structure are
  $A_1(q), A_2(q), . . . , A_t(q)$, return the median $M(q)$ of these
  $t$ answers.  Show that with $t = c_1 \log \frac{1+}{\delta}$, and
  $d = \frac{c_2}{\epsilon^2}$, where $c_1, c_2$ are constants you can
  choose, you get $M(q) \in x_q \pm \epsilon || x ||_2$ with
  probability $1-\delta$.
\end{enumerate}

\litem{(Chernoff meets Matrices)} In Lecture 18 we mentioned a very
general theorem about matrix-valued Chernoff bounds for symmetric
matrices. In this problem we’ll take the first steps towards
it. Assume eigenvalues are numbered so that
$\lambda_1 \geq \ldots \lambda_n$. We’ll prove:
\begin{theorem}
  Let $X_1, X_2, \ldots, X_n$ be independent symmetric $d \times t$ matrices, and $S_n= \sum_i X_)i$. Then, for any $t \geq 0$ and $\ell \in \R$,
  \[\P[\lambda_1(S_n) \geq \ell] \leq d \cdot e^{-t\ell} \cdot \prod_{i=1}^n \lambda_1 (\E[e^{tX_i}]).\]
    \[\P[\lambda_d(S_n) \leq -\ell] \leq d \cdot e^{-t\ell} \cdot \prod_{i=1}^n \lambda_1 (\E[e^{-tX_i}]).\]
  \end{theorem}
  Recall that $\tr(A) = \sum_{i=1}^n a_{ii}$. You may use the following without proof:
\begin{enumerate}

  \begin{enumerate}
  \item $\tr(A) = \sum_i \lambda_i (A)$
  \item $\lambda_i (e^A) = e^{\lambda_i(A)}$
  \item The Golden-Thompson inequality: $\tr(e^{A+B}) \leq \tr(e^A \cdot e^B)$.
  \item For PSD matrices $A$ and $B$, $\tr(AB) \leq \tr(A)\cdot \tr(B)$.\
  \item Expectations and trace commute: $\E[\tr(A)] = \tr(\E[A])$.
  \end{enumerate}
\item Show that for any $t \geq 0$,
  \[\P[\lambda_1(S_n) \geq \ell] \leq \P[\tr(e^{tS_n}) \geq e^{t\ell}] \leq e^{-t\ell} \cdot \E[\tr(s^{tS_n}].\]
\item Show that
  \[\E_{X_1, \ldots, X_n}[\tr(e^{tS_n})]\leq \E_{X_1, \ldots, X_{n-1}} [\tr(e^{tS_{n-1}})]\cdot \lambda_1 (\E[e^{tX_n}]).\]
  Hint: why can you use $(iv)$ above even if $X_n$ isn't PSD?
\item Use the previous two parts to prove the first statement of the theorem.
  \item Use the same arguments on $(-S_n) = \sum_i (-X_i)$ to prove the other part.
\end{enumerate}
\end{enumerate}




\end{document}
